\scanpage{78}\mspage{65}
% archive sheet number: 77

3.10. Ainsi, pour avoir un dérivateur au plein
sens du terme, soit sur $\mathrm{Diag}_0$, soit sur
$\mathrm{Diag}$, il ne manque que Der 4 (existence
des $u_*$), Der 5 (calcul des $u_*$) et
Der 6 (exactitude des carrés $W$\deleted{\unsure{co}}-cartésiens).
On va \deleted{\unsure{l'}}\inserted{les} obtenir en principe à l'aide
d'une structure de \uline{catégorie à fibrations}
$(M,W,F)$ associée à $(M,W)$, \deleted{\unsure{ce}} i.e.
l'espèce de structure duale de celle
introduite \deleted{\unsure{dans}} p. 39. Ces \inserted{deux} structures
n'ont pas à être reliées entre elles d'une
façon \unsure{étroite|autre} que par le fait qu'elles
correspondent au $\unsure{=}$ $W$. Elles sont
considérablement plus faciles à \deleted{\unsure{dégag}}
construire que des structures de Quillen
$(W,C,F)$, où $C,F$ doivent être reliés
de façon \unsure{autrement|extrêmement} stricte, \deleted{\unsure{savoir}} \unsure{aux|par}
\deleted{\unsure{des prop}} les propriétés $C\longleftrightarrow\overset{F\cap W}{\overset{\shortparallel}{TF}}$, $\overset{C\cap W}{\overset{\shortparallel}{TC}}\unsure{\longleftarrow|\longleftrightarrow} F$,
(dont il n'est pas question dans l'axiomatique
de K.S. Brown), et \deleted{\unsure{de}}factorisation pour
$\unsure{\text{--}|(}C,TF)$ et $(TC,F)$.
